\chapter{Conclusions and Future Work} 
This chapter concludes the thesis, evaluating how well the goals we set in the introduction were reached. We look back at the results and evaluate the work done so far. In the final section we look at future work to be done on this project. 


\section{Conclusions}
In Chapter~\ref{chap:introduction} of this thesis we posed two questions: 
\begin{enumerate}
	\item Can a partition-aware objective produce better balanced communication loads than a global one?
	\item Can the performance of a parallel application be improved in practice by using a partition-aware objective?
\end{enumerate}

\subsection{Contributions}
%what was built 
%bullet list of contributions? idk about this one
%set up next subsections
WIP NOT FINISHED

To answer them, we developed NVOL, a novel implementation of a partition-aware objective function based on METIS. Chapters~\ref{cv_calculation} and~\ref{chapter5} describe how the objective function was developed, and Chapter~\ref{chap:NVOL tooling} describes how it was integrated into METIS.

\subsection{Question 1 - Does NVOL improve the volume balance?}
Our results show that NVOL generates partitions with volumes comparable to those produced by METIS when optimizing for total communication volume. NVOL focuses on minimizing the maximum outgoing volume ($W^{\mathrm{max}}$) of any single partition, allowing moves that raise the volume elsewhere as long as the maximum is not increased. This consistently yields partitions with a lower maximum and higher minimum outgoing volume compared to VOL (METIS's total communication volume objective function). The more even distribution is reflected in the lower variance in outgoing volumes, especially in runs with 2 partitions, where NVOL produces perfectly balanced partitions on almost all graphs.

On more regular graphs, NVOL consistently improves on $W^{\mathrm{max}}$. The numerical simulation graphs and the cardiac graphs are particularly notable, with NVOL achieving a lower $W^{\mathrm{max}}$ than VOL in almost every run. On these graphs NVOL also reduces the spread of outgoing volumes across partitions, producing more uniform partition volumes. Notably, on the cardiac graphs and a few others, NVOL even achieves a lower total communication volume than VOL despite not optimizing for that metric.

These results are not universal: NVOL clearly struggles on certain graph structures. Road network graphs in particular proved challenging in our tests, this was the only graph type where NVOL failed to improve on $W^{\mathrm{max}}$ in as many as 45\% of runs. These graphs are sparse, with a few high-degree vertices and many low-degree ones, and this irregular structure prevents NVOL from producing consistent improvements.

The largest graph we managed to test NVOL on was kmer\_U1a, a protein k-mer graph with 67.7 million vertices. This graph is also sparse, with many small, highly connected components. NVOL nevertheless produces a better $W^{\mathrm{max}}$ than VOL and consistently lower standard deviations in the volume values, reaching a reduction in $W^{\mathrm{max}}$ of up to 46.34\%. This suggests that NVOL performs well when given many options, and raises an interesting question: how well would NVOL perform on even larger graphs? On the more negative side, this graph was already too large for NVOL to partition into 2048 parts, making efficiency improvements a priority before scaling further.

This point about NVOL performing better when it has many options is also supported by the results for coPapersDBLP, a graph with a wide variance in edges per vertex. This structure is common in citation networks, which have a few high-degree hub vertices and many low-degree ones. NVOL struggles considerably at low $k$, with the worst result at $-32.09\%$. The trend reverses as $k$ grows, however, peaking at an improvement of $56.70\%$.

We also tested NVOL with imbalance tuning, which produced more variable results. Allowing greater imbalance introduced more instability, with several additional runs failing to improve on $W^{\mathrm{max}}$, though NVOL maintained a positive win rate overall. The main effect of relaxed imbalance appears to be a larger magnitude in both wins and losses. The high variance of these results suggests that imbalance tuning should be treated as a per-run option rather than a default setting.

Taken together, these results answer the first question affirmatively: NVOL produces partitions with more uniform volumes and lower variance compared to METIS's volume objective function. This does push the average partition volume higher, but the increase comes primarily from redistributing the spikes that previously concentrated on a few partitions. NVOL improves the maximum outgoing volume in the partitions, achieving results up to $56.70\%$ compared to METIS's total volume objective.

\subsection{Question 2 - Do the partitions produced by NVOL improve performance?}
The answer to this question is yes, but with an asterisk. NVOL achieves better runtimes than VOL in 65\% of our tests, and maintains a 60\% win rate when imbalance tuning is enabled. On its wins, NVOL improves runtime by an average of 2.61\%, peaking at 6.6\% (without imbalance tuning). 

Some results, such as those for italy\_osm, also show that latency can dominate volume in certain settings. The time spent transmitting the data itself is negligible compared to the latency of initiating each message. As a result, on italy\_osm the tuned NVOL produced a runtime 27.8\% slower than VOL, even though it achieved a 26.61\% lower $W^{\mathrm{max}}$. The slowdown came purely from an increase in the number of partition neighbors: with more partitions to message, the per-message latency dominated the overall runtime.

The coPapersDBLP graph highlights another issue: by producing a higher average volume, NVOL slows down the more heavily loaded processes. Other processes must then wait for them to finish communicating, so the increased per-partition volume affects the overall runtime negatively. This suggests that NVOL will struggle on irregular graphs such as road and citation networks, particularly when they are partitioned into a small number of parts.

These results come with important caveats. We were only able to test on an idealized setup with 32 processors, which significantly limited the number of configurations we could explore. Most of NVOL's greatest improvements appear on partitionings into a large number of parts. Furthermore, these values are only estimates, providing no definitive confirmation that NVOL's runtimes are actually faster on real hardware.

In summary, NVOL does appear to improve efficiency in many cases, but these improvements remain simulated. Future work would require additional testing on real machines to validate these findings. Even so, the estimates we do have are promising and suggest that NVOL's partitions could lead to meaningful performance improvements, especially on graphs partitioned into many parts.

\section{Future Work}\label{chap_future_work}
NVOL shows good potential, delivering consistently strong results across a wide variety of graphs. The conclusions above point out two natural directions for future work: reducing the partitioning time overhead introduced by NVOL, and extending the objective function with hardware-aware information to better target real parallel machines.

\subsection{Improved efficiency} \label{sec:Improved eff}
One of METIS's main strengths is its partitioning speed, and the additional structures and operations introduced by NVOL inevitably affect this, particularly on larger graphs. For small graphs the partitioning time is essentially unchanged, but on larger ones NVOL was observed to be significantly slower than VOL.

The overhead has a straightforward explanation: NVOL adds code to METIS's refinement stage and increases the number of iterations performed there. METIS itself is carefully designed to avoid redundant work, most of the volume operations are performed during graph projection. This way the refinement stage can reuse it rather than reiterating over the graph. NVOL breaks this pattern by recomputing the gain information for each vertex it visits during refinement, effectively repeating work that METIS already performs at projection time.

Three considerations led us not to integrate NVOL's gain calculation directly into the projection stage. First, NVOL has a larger memory requirement than VOL. As described in Chapter~\ref{cv_calculation}, VOL requires a vector whose length equals the number of partitions a vertex neighbors, whereas NVOL requires a table. For a vertex with a gain vector of length $k$, the corresponding gain table has dimensions $(k+1)\times k$, so a gain vector of length 5 already implies a $6 \times 5$ table per vertex. Second, VOL uses incremental updates between moves to keep the volume gains of affected vertices current; implementing the equivalent logic for NVOL was outside the scope of this thesis. Third, and perhaps most decisive, time constraints during development, combined with the considerable effort required to address the first two points, prevented us from pursuing this integration further. 

This decision was nonetheless preceded by an attempt at the integration. Our strategy allocated the exact memory needed for each gain table and deallocated it once the corresponding vertex left the boundary; gain updates were handled by simply recomputing all neighbors of the moved vertex. In principle this is much closer to what VOL does and should have improved performance. In practice, however, the algorithm became even slower than the original NVOL implementation, with the frequent allocation and deallocation of memory as the dominant cost. A better solution would be to allocate a single block of memory at projection time, sized from an estimate of how many boundary vertices the finer graph will contain, with additional logic to extend the block when the estimate falls short and to release it once refinement completes. Combined with proper incremental updates of neighboring vertices' gains after each move, this would allow NVOL to partition in a time frame close to VOL, at the cost of higher memory use.

\subsection{Introduce hardware knowledge}
Using the partition-specific implementation of NVOL, it would be interesting to incorporate some of the logic presented in~\cite{Xing_of_Heterogeneous} to further improve actual runtime. Currently, we partition using communication volume and then model the runtime of each partition based on a set of constants and the incoming volume. With the right adaptations, NVOL could already evaluate this during refinement. Changing NVOL's volume calculations from outgoing to incoming volume is non-trivial, but it has the potential for even larger improvements.

The general idea would be to extend METIS to accept additional parameters describing the latencies and bandwidth of the machine running the partitions. Then, during refinement, one could use the equations shown in Section~\ref{sec:model estimates} to evaluate candidate moves. This would increase the chance of producing a partitioning that performs well on the target parallel machine. NVOL's per-partition structure is essential to such an approach.

This could be extended further to optimize for heterogeneous machines. So far we have only considered tests on machines where each component was identical, but in reality modern parallel computers are larger and more complex. Their heterogeneous structure therefore brings an added level of complexity that can be exploited. In~\cite{Xing_of_Heterogeneous}, the authors also develop models for such setups. With the right structures in place, NVOL could support this logic as well.